% ======================================================================
% Chapter 9: Functional Methods
% (book pages 275-314; PDF pages 298-337)
% ======================================================================
\chapter{Functional Methods}

Feynman once said that%
\footnote{\emph{The Character of Physical Law} (MIT Press, 1965), p.~168.}
``every theoretical physicist who is any good knows six or seven different
theoretical representations for exactly the same physics.''  Following his
advice, we introduce in this chapter an alternative method of deriving the
Feynman rules for an interacting quantum field theory: the method of
\emph{functional integration}.

Aside from Feynman's general principle, we have several specific reasons for
introducing this formalism.  It will provide us with a relatively easy
derivation of our expression for the photon propagator, completing the proof
of the Feynman rules for QED given in Section~4.8.  The functional method
generalizes more readily to other interacting theories, such as scalar QED
(Problem~9.1), and especially the non-Abelian gauge theories (Part~III).
Since it uses the Lagrangian, rather than the Hamiltonian, as its fundamental
quantity, the functional formalism explicitly preserves all symmetries of a
theory.  Finally, the functional approach reveals the close analogy between
quantum field theory and statistical mechanics.  Exploiting this analogy, we
will turn Feynman's advice upside down and apply the same theoretical
representation to two completely different areas of physics.

\section{Path Integrals in Quantum Mechanics}

We begin by applying the functional integral (or path integral) method to the
simplest imaginable system: a nonrelativistic quantum-mechanical particle
moving in one dimension.  The Hamiltonian for this system is
\begin{equation}
  H = \frac{p^2}{2m} + V(q) .
\end{equation}
Suppose that we wish to compute the amplitude for this particle to travel from
one point ($x_a$) to another ($x_b$) in a given time ($T$).  We will call this
amplitude $U(x_a, x_b; T)$; it is the position representation of the
Schr\"odinger time-evolution operator.  In the canonical Hamiltonian
formalism, $U$ is given by
\begin{equation}
  U(x_a, x_b; T) = \langle x_b|\ e^{-iHT/\hbar}\ |x_a\rangle .
  \label{eq:9.1}
\end{equation}
(For the next few pages we will display all factors of $\hbar$ explicitly.)

In the path-integral formalism, $U$ is given by a very different-looking
expression.  We will first try to motivate that expression, then prove that it
is equivalent to \eqref{eq:9.1}.

Recall that in quantum mechanics there is a superposition principle: When a
process can take place in more than one way, its total amplitude is the
coherent sum of the amplitudes for each way.  A simple but nontrivial example
is the famous double-slit experiment, shown in Fig.~\ref{fig:9.1}.  The total
amplitude for an electron to arrive at the detector is the sum of the
amplitudes for the two paths shown.  Since the paths differ in length, these
two amplitudes generally differ, causing interference.

\begin{figure}[t]
  \centering
  \includegraphics[width=0.55\textwidth]{fig9.1.png}
  \caption{The double-slit experiment.  Path 2 is longer than path 1 by an
    amount $d$, and therefore has a phase that is larger by
    $2\pi d/\lambda$, where $\lambda = 2\pi\hbar/p$ is the particle's de
    Broglie wavelength.  Constructive interference occurs when
    $d = 0, \lambda, \ldots$, while destructive interference occurs when
    $d = \lambda/2, 3\lambda/2, \ldots$.}
  \label{fig:9.1}
\end{figure}

For a general system, we might therefore write the total amplitude for
traveling from $x_a$ to $x_b$ as
\begin{equation}
  \langle x_b|e^{-iHT/\hbar}|x_a\rangle = \sum_{\text{all paths}}
    \text{(phase)} .
  \label{eq:9.2}
\end{equation}
To be democratic, we have written the amplitude for each particular path as a
pure phase, so that no path is inherently more important than any other.  The
symbol $\int\mathcal{D}x(t)$ is simply another way of writing ``sum over all
paths''; since there is one path for every function $x(t)$ that begins at
$x_a$ and ends at $x_b$, the sum is actually an integral over this continuous
space of functions.

We can define this integral as part of a natural generalization of the
calculus to spaces of functions.  A function that maps functions to numbers is
called a \emph{functional}.  The integrand in \eqref{eq:9.2} is a functional,
since it associates a complex amplitude with any function $x(t)$.  The
argument of a functional $F[x(t)]$ is conventionally written in square
brackets rather than parentheses.  Just as an ordinary function $y(x)$ can be
integrated over a set of points $x$, a functional $F[x(t)]$ can be integrated
over a set of functions $x(t)$; the measure of such a functional integral is
conventionally written with a script capital $D$, as in \eqref{eq:9.2}.  A
functional can also be differentiated with respect to its argument (a
function), and this \emph{functional derivative} is denoted by
$\delta F/\delta x(t)$.  We will develop more precise definitions of this new
integral and derivative in the course of this section and the next.

What should we use for the ``phase'' in Eq.~\eqref{eq:9.2}?  In the classical
limit, we should find that only one path, the classical path, contributes to
the total amplitude.  We might therefore hope to evaluate the integral in
\eqref{eq:9.2} by the method of stationary phase, identifying the classical
path $x_{cl}(t)$ by the stationary condition,
\begin{equation}
  \frac{\delta S[x]}{\delta x(t)}\Big|_{x=x_{cl}} = 0 .
\end{equation}
But the classical path is the one that satisfies the principle of least
action,
\begin{equation}
  \delta S = 0,\qquad S = \int L\,dt ,
\end{equation}
where $S = \int L\,dt$ is the classical action.  It is tempting, therefore, to
identify the phase with $S$, up to a constant.  Since the stationary-phase
approximation should be valid in the classical limit---that is, when
$S \gg \hbar$---we will use $S/\hbar$ for the phase.  Our final formula for the
propagation amplitude is thus
\begin{equation}
  \langle x_b|\ e^{-iHT/\hbar}\ |x_a\rangle = U(x_a, x_b; T)
  = \int \mathcal{D}x(t)\ e^{iS[x(t)]/\hbar} .
  \label{eq:9.3}
\end{equation}
We can easily verify that this formula gives the correct interference pattern
in the double-slit experiment.  The action for either path shown in
Fig.~\ref{fig:9.1} is just $(1/2)mv^2t$, the kinetic energy times the time.
For path 1 the velocity is $v_1 = D/t$, so the phase is $mD^2/2\hbar t$.  For
path 2 we have $v_2 = (D+d)/t$, so the phase is $m(D+d)^2/2\hbar t$.  We must
assume that $d \ll D$, so that $v_1 \approx v_2$ (i.e., the electrons have a
well-defined velocity).  The excess phase for path 2 is then
$mDd/\hbar t \approx pd/\hbar$, where $p$ is the momentum.  This is exactly
what we would expect from the de Broglie relation $p = \hbar/\lambda$, so we
must be doing something right.

To evaluate the functional integral more generally, we must define the symbol
$\int\mathcal{D}x(t)$ in the case where the number of paths $x(t)$ is more
than two (and, in fact, continuously infinite).  We will use a brute-force
definition, by discretization.  Break up the time interval from 0 to $T$ into
many small pieces of duration $\epsilon$, as shown in Fig.~\ref{fig:9.2}.
Approximate a path $x(t)$ as a sequence of straight lines, one in each time
slice.  The action for this discretized path is
\begin{equation}
  \sum_k \epsilon\ \Bigl[ \frac{m}{2}\Bigl( \frac{x_{k+1}-x_k}{\epsilon}
    \Bigr)^2 - V\Bigl( \frac{x_{k+1}+x_k}{2} \Bigr) \Bigr] .
\end{equation}

\begin{figure}[t]
  \centering
  \includegraphics[width=0.5\textwidth]{fig9.2.png}
  \caption{We define the path integral by dividing the time interval into
    small slices of duration $\epsilon$, then integrating over the coordinate
    $x_k$ of each slice.}
  \label{fig:9.2}
\end{figure}

We then define the path integral by
\begin{equation}
  \begin{aligned}
    \int \mathcal{D}x(t)\ F[x(t)] &= \lim_{\epsilon\to 0}\
      \Bigl[ \prod_k \int \frac{dx_k}{C(\epsilon)} \Bigr]\ F[x(t)]\,
      e^{i\sum_k \epsilon \mathcal{L}_k}\\
    &= \lim_{\epsilon\to 0}\ \Bigl[ \int \frac{dx_1}{C(\epsilon)}
      \cdots \int \frac{dx_{N-1}}{C(\epsilon)} \Bigr]\ F[x(t)]\,
      e^{i\sum_k \epsilon \mathcal{L}_k} ,
  \end{aligned}
  \label{eq:9.4}
\end{equation}
where $C(\epsilon)$ is a constant, to be determined later.  (We have included
one factor of $C(\epsilon)$ for each of the $N$ time slices, for reasons that
will be clear below.)  At the end of the calculation we take the limit
$\epsilon\to 0$.  (As in Sections~4.5 and~6.2, the $\prod$ symbol is an
instruction to write what follows once for each $k$.)

Using \eqref{eq:9.4} as the definition of the right-hand side of
\eqref{eq:9.3}, we will now demonstrate the validity of \eqref{eq:9.3} for a
general one-particle potential problem.  To do this, we will show that the
left- and right-hand sides of \eqref{eq:9.3} are obtained by integrating the
same differential equation, with the same initial condition.  In the process,
we will determine the constant $C(\epsilon)$.

To derive the differential equation satisfied by \eqref{eq:9.4}, consider the
addition of the very last time slice in Fig.~\ref{fig:9.2}.  According to
\eqref{eq:9.3} and the definition \eqref{eq:9.4}, we should have
\begin{equation}
  \begin{aligned}
    U(x_a, x_b; T) &= \int_{-\infty}^{\infty} dx'\ C(\epsilon)\
      \exp\Bigl[ \frac{im(x_b-x')^2}{2\hbar\epsilon}
        - \frac{i\epsilon}{\hbar}V\Bigl( \frac{x_b+x'}{2} \Bigr) \Bigr]\
      U(x_a, x'; T-\epsilon) .
  \end{aligned}
\end{equation}
The integral over $x'$ is just the contribution to $\int\mathcal{D}x$ from the
last time slice, while the exponential factor is the contribution to
$e^{iS/\hbar}$ from that slice.  All contributions from previous slices are
contained in $U(x_a, x'; T-\epsilon)$.

As we send $\epsilon\to 0$, the rapid oscillation of the first term in the
exponential constrains $x'$ to be very close to $x_b$.  We can therefore
expand the above expression in powers of $(x'-x_b)$:
\begin{equation}
  U(x_a, x_b; T) = \int_{-\infty}^{\infty} dx'\ C(\epsilon)\
    \exp\Bigl[ \frac{im(x_b-x')^2}{2\hbar\epsilon} \Bigr]\
    \Bigl[ 1 - \frac{i\epsilon}{\hbar}V(x_b) + \cdots \Bigr]\
    U(x_a, x_b; T-\epsilon) .
  \label{eq:9.5}
\end{equation}
We can now perform the $x'$ integral by treating the exponential factor as a
Gaussian.  (Properly, we should introduce a small real term in the exponent
for convergence; we will ignore this term until the next section, when we
derive Feynman rules using functional methods.)  Recall the Gaussian
integration formulae
\begin{equation}
  \int d\xi\ e^{i b\xi^2} = \sqrt{\frac{i\pi}{b}},\qquad
  \int d\xi\ \xi\,e^{i b\xi^2} = 0,\qquad
  \int d\xi\ \xi^2 e^{i b\xi^2} = \frac{i}{2b}\sqrt{\frac{i\pi}{b}} .
\end{equation}
Applying these identities to \eqref{eq:9.5}, we find
\begin{equation}
  U(x_a, x_b; T) = C(\epsilon)\sqrt{\frac{2\pi i\hbar\epsilon}{-m}}\
    \Bigl( 1 - \frac{i\epsilon}{\hbar}V(x_b) + \frac{i\hbar\epsilon}{2m}
      \frac{\partial^2}{\partial x_b^2} + \mathcal{O}(\epsilon^2) \Bigr)\
    U(x_a, x_b; T-\epsilon) .
\end{equation}
This expression makes no sense in the limit $\epsilon\to 0$ unless the factor
in parentheses is equal to 1.  We can therefore identify the correct definition
of $C$:
\begin{equation}
  C(\epsilon) = \Bigl( \frac{2\pi\hbar\epsilon}{-im} \Bigr)^{1/2}
  = \sqrt{\frac{2\pi i\hbar\epsilon}{m}} .
  \label{eq:9.6}
\end{equation}
Given this definition, we can compare terms of order $\epsilon$ and multiply
by $i\hbar$ to obtain
\begin{equation}
  i\hbar\frac{\partial}{\partial T}U(x_a, x_b; T)
  = \Bigl[ -\frac{\hbar^2}{2m}\frac{\partial^2}{\partial x_b^2}
    + V(x_b) \Bigr]U(x_a, x_b; T)
  = H U(x_a, x_b; T) .
  \label{eq:9.7}
\end{equation}
This is the Schr\"odinger equation.  But it is easy to show that the
time-evolution operator $U$, as originally defined in \eqref{eq:9.1},
satisfies the same equation.

As $T\to 0$, the left-hand side of \eqref{eq:9.3} tends to
$\delta(x_a-x_b)$.  Compare this to the value of \eqref{eq:9.4} in the case of
one time slice:
\begin{equation}
  \frac{1}{C(\epsilon)}\ \exp\Bigl[ \frac{im(x_b-x_a)^2}{2\hbar\epsilon}
    \Bigr] \to \delta(x_a-x_b) + \mathcal{O}(\epsilon) .
\end{equation}
This is just the peaked exponential of \eqref{eq:9.5}, and it also tends to
$\delta(x_a-x_b)$ as $\epsilon\to 0$.  Thus the left- and right-hand sides of
\eqref{eq:9.3} satisfy the same differential equation with the same initial
condition.  We conclude that the Hamiltonian definition of the time evolution
operator \eqref{eq:9.1} and the path-integral definition \eqref{eq:9.3} are
equivalent, at least for the case of this simple one-dimensional system.

To conclude this section, let us generalize our path-integral formula to more
complicated quantum systems.  Consider a very general quantum system,
described by an arbitrary set of coordinates $q^i$, conjugate momenta $p_i$,
and Hamiltonian $H(q,p)$.  We will give a direct proof of the path-integral
formula for transition amplitudes in this system.

The transition amplitude that we would like to compute is
\begin{equation}
  U(q_a, q_b; T) = \langle q_b|\ e^{-iHT}\ |q_a\rangle .
  \label{eq:9.8}
\end{equation}
(When $q$ or $p$ appears without a superscript, it will denote the set of all
coordinates $\{q^i\}$ or momenta $\{p_i\}$.  Also, for convenience, we now set
$\hbar = 1$.)  We break up the time interval into $N$ slices of duration
$\epsilon = T/N$, and insert complete sets of position and momentum states
between successive time slices.  The result is
\begin{equation}
  \begin{aligned}
    U(q_a, q_b; T) &= \Bigl[ \prod_{k=1}^{N-1}\int dq_k \Bigr]\
      \Bigl[ \prod_{k=0}^{N-1}\int \frac{dp_k}{2\pi} \Bigr]\
      \exp\Bigl\{ i\sum_{k=0}^{N-1}\Bigl[ p_k(q_{k+1}-q_k)
        - \epsilon H(p_k, \bar{q}_k) \Bigr] \Bigr\}\\
    &= \int \mathcal{D}q(t)\,\mathcal{D}p(t)\ \exp\Bigl[ i\int_0^T dt\
      \bigl( p\dot{q} - H(p,q) \bigr) \Bigr] .
  \end{aligned}
  \label{eq:9.12}
\end{equation}
This is the phase-space form of the path integral.

For a nonrelativistic particle, the Hamiltonian is simply
$H = p^2/2m + V(q)$.  In this case we can evaluate the $p$-integrals by
completing the square in the exponent:
\begin{equation}
  \int \frac{dp_k}{2\pi}\ \exp\Bigl[ i\Bigl( p_k(q_{k+1}-q_k)
    - \epsilon\frac{p_k^2}{2m} \Bigr) \Bigr]
  = C(\epsilon)\ \exp\Bigl[ \frac{im}{2\epsilon}(q_{k+1}-q_k)^2 \Bigr] ,
\end{equation}
where $C(\epsilon)$ is just the factor \eqref{eq:9.6}.  Notice that we have one
such factor for each time slice.  Thus we recover expression \eqref{eq:9.3},
in discretized form, including the proper factors of $C$:
\begin{equation}
  U(q_a, q_b; T) = \lim_{\epsilon\to 0}\ \Bigl[ \prod_k
    \int \frac{dq_k}{C(\epsilon)} \Bigr]\
    \exp\Bigl\{ \sum_k \epsilon\ \Bigl[ \frac{m}{2}\Bigl(
      \frac{q_{k+1}-q_k}{\epsilon} \Bigr)^2 - V(q_k) \Bigr] \Bigr\}.
  \label{eq:9.13}
\end{equation}

\section{Functional Quantization of Scalar Fields}

In this section we will apply the functional integral formalism to the quantum
theory of a real scalar field $\phi(x)$.  Our goal is to derive the Feynman
rules for such a theory directly from functional integral expressions.

The general functional integral formula \eqref{eq:9.12} derived in the last
section holds for any quantum system, so it should hold for a quantum field
theory.  In the case of a real scalar field, the coordinates $q^i$ are the
field amplitudes $\phi(\mathbf{x})$, and the Hamiltonian is
\begin{equation}
  H = \int d^3x\ \Bigl[ \tfrac{1}{2}\pi^2 + \tfrac{1}{2}(\nabla\phi)^2
    + V(\phi) \Bigr].
\end{equation}
Thus our formula becomes
\begin{equation}
  \langle \phi_b(\mathbf{x})|\ e^{-iHT}\ |\phi_a(\mathbf{x})\rangle
  = \int \mathcal{D}\phi\,\mathcal{D}\pi\ \exp\Bigl[ i\int_0^T d^4x\
    \bigl( \pi\dot{\phi} - \tfrac{1}{2}\pi^2 - \tfrac{1}{2}(\nabla\phi)^2
      - V(\phi) \bigr) \Bigr] ,
\end{equation}
where the functions $\phi(\mathbf{x})$ over which we integrate are constrained
to the specific configurations $\phi_a(\mathbf{x})$ at $x^0 = 0$ and
$\phi_b(\mathbf{x})$ at $x^0 = T$.  Since the exponent is quadratic in $\pi$,
we can complete the square and evaluate the $\int\mathcal{D}\pi$ integral to
obtain
\begin{equation}
  \langle \phi_b(\mathbf{x})|\ e^{-iHT}\ |\phi_a(\mathbf{x})\rangle
  = \int_{\phi_a}^{\phi_b} \mathcal{D}\phi\ \exp\Bigl[ i\int_0^T d^4x\
    \mathcal{L} \Bigr],
  \label{eq:9.14}
\end{equation}
where
\begin{equation}
  \mathcal{L} = \tfrac{1}{2}(\partial_\mu\phi)^2 - V(\phi)
\end{equation}
is the Lagrangian density.  The integration measure $\mathcal{D}\phi$ in
\eqref{eq:9.14} again involves an awkward constant, which we will not write
explicitly.

The time integral in the exponent of \eqref{eq:9.14} goes from 0 to $T$, as
determined by our choice of what transition function to compute; in all other
respects this formula is manifestly Lorentz invariant.  Any other symmetries
that the Lagrangian may have are also explicitly preserved by the functional
integral.  As we proceed in our study of quantum field theory, symmetries and
their associated conservation laws will play an increasingly central role.  We
therefore propose to take a rash step: Abandon the Hamiltonian formalism, and
take Eq.~\eqref{eq:9.14} to define the Hamiltonian dynamics.  Any such formula
corresponds to some Hamiltonian; to find it, one can always differentiate with
respect to $T$ and derive the Schr\"odinger equation as in the previous
section.  We thus consider the Lagrangian $\mathcal{L}$ to be the most
fundamental specification of a quantum field theory.  We will see next that
one can use the functional integral to compute from $\mathcal{L}$ directly,
without invoking the Hamiltonian at all.

\subsection{Correlation Functions}

To make direct use of the functional integral, we need a functional formula
for computing correlation functions.  To find such an expression, consider the
object
\begin{equation}
  \int \mathcal{D}\phi(x)\ \phi(x_1)\phi(x_2)\ \exp\Bigl[ i\int_{-T}^{T}
    d^4x\ \mathcal{L}(\phi) \Bigr],
  \label{eq:9.15}
\end{equation}
where the boundary conditions on the path integral are
$\phi(-T,\mathbf{x}) = \phi_a(\mathbf{x})$ and
$\phi(T,\mathbf{x}) = \phi_b(\mathbf{x})$ for some $\phi_a$, $\phi_b$.  We
would like to relate this quantity to the two-point correlation function,
$\langle\Omega|T\phi_H(x_1)\phi_H(x_2)|\Omega\rangle$.  (To distinguish
operators from ordinary numbers, we write the Heisenberg picture operator
with an explicit subscript: $\phi_H(x)$.  Similarly, we will write
$\phi_S(\mathbf{x})$ for the Schr\"odinger picture operator.)

First we break up the functional integral in \eqref{eq:9.15} as follows:
\begin{equation}
  \int \mathcal{D}\phi(x) = \int \mathcal{D}\phi_1(\mathbf{x})\,
    \int \mathcal{D}\phi_2(\mathbf{x})\, \int_{\phi_1}^{\phi_2}
    \mathcal{D}\phi(x) .
  \label{eq:9.16}
\end{equation}
The main functional integral $\int\mathcal{D}\phi(x)$ is now constrained at
times $x^0_1$ and $x^0_2$ (in addition to the endpoints $-T$ and $T$), but we
must integrate separately over the intermediate configurations
$\phi_1(\mathbf{x})$ and $\phi_2(\mathbf{x})$.  After this decomposition, the
extra factors $\phi(x_1)$ and $\phi(x_2)$ in \eqref{eq:9.15} become
$\phi_1(x_1)$ and $\phi_2(x_2)$, and can be taken outside the main integral.
The main integral then factors into three pieces, each being a simple
transition amplitude according to \eqref{eq:9.14}.  The times $x^0_1$ and
$x^0_2$ automatically fall in order; for example, if $x^0_1 < x^0_2$, then
\eqref{eq:9.15} becomes
\begin{equation}
  \int \mathcal{D}\phi_1(\mathbf{x})\int \mathcal{D}\phi_2(\mathbf{x})\
    \phi_1(x_1)\phi_2(x_2)\ \langle \phi_b|e^{-iH(T-x^0_2)}|\phi_2\rangle\
    \langle \phi_2|e^{-iH(x^0_2-x^0_1)}|\phi_1\rangle\
    \langle \phi_1|e^{-iH(x^0_1+T)}|\phi_a\rangle .
\end{equation}
We can turn the field $\phi_1(x_1)$ into a Schr\"odinger operator using
$\phi_S(\mathbf{x})|\phi_1\rangle = \phi_1(\mathbf{x})|\phi_1\rangle$.  The
completeness relation $\int\mathcal{D}\phi_1\,|\phi_1\rangle\langle\phi_1| = 1$
then allows us to eliminate the intermediate state $|\phi_1\rangle$.  Similar
manipulations work for $\phi_2$, yielding the expression
\begin{equation}
  \langle \phi_b|\ e^{-iH(T-x^0_2)}\ \phi_S(\mathbf{x}_2)\
    e^{-iH(x^0_2-x^0_1)}\ \phi_S(\mathbf{x}_1)\
    e^{-iH(x^0_1+T)}\ |\phi_a\rangle .
\end{equation}
Most of the exponential factors combine with the Schr\"odinger operators to
make Heisenberg operators.  In the case $x^0_1 > x^0_2$, the order of $x_1$ and
$x_2$ would simply be interchanged.  Thus expression \eqref{eq:9.15} is equal
to
\begin{equation}
  \langle \phi_b|\ e^{-iHT}\ T\{\phi_H(x_1)\phi_H(x_2)\}\ e^{-iHT}\
    |\phi_a\rangle .
  \label{eq:9.17}
\end{equation}
This expression is almost equal to the two-point correlation function.  To
make it more nearly equal, we take the limit $T\to\infty(1-i\epsilon)$.  Just
as in Section~4.2, this trick projects out the vacuum state $|\Omega\rangle$
from $|\phi_a\rangle$ and $|\phi_b\rangle$ (provided that these states have
some overlap with $|\Omega\rangle$, which we assume).  For example, decomposing
$|\phi_a\rangle$ into eigenstates $|n\rangle$ of $H$, we have
\begin{equation}
  e^{-iHT}|\phi_a\rangle = \sum_n e^{-iE_nT}|n\rangle\langle n|\phi_a\rangle
  \xrightarrow{T\to\infty(1-i\epsilon)}
    \langle\Omega|\phi_a\rangle\ e^{-iE_0T}|\Omega\rangle .
\end{equation}
As in Section~4.2, we obtain some awkward phase and overlap factors.  But
these factors cancel if we divide by the same quantity as \eqref{eq:9.15} but
without the two extra fields $\phi(x_1)$ and $\phi(x_2)$.  Thus we obtain the
simple formula
\begin{equation}
  \langle\Omega|T\phi_H(x_1)\phi_H(x_2)|\Omega\rangle
  = \lim_{T\to\infty(1-i\epsilon)}\
    \frac{\int\mathcal{D}\phi\ \phi(x_1)\phi(x_2)\ \exp\bigl[
      i\int_{-T}^{T}d^4x\ \mathcal{L} \bigr]}
      {\int\mathcal{D}\phi\ \exp\bigl[ i\int_{-T}^{T}d^4x\ \mathcal{L} \bigr]} .
  \label{eq:9.18}
\end{equation}
This is our desired formula for the two-point correlation function in terms of
functional integrals.  For higher correlation functions, just insert additional
factors of $\phi$ on both sides.

\subsection{Feynman Rules}

Our next task is to compute various correlation functions directly from the
right-hand side of formula \eqref{eq:9.18}.  In other words, we will now use
\eqref{eq:9.18} to derive the Feynman rules for a scalar field theory.  We
will begin by computing the two-point function in the free Klein--Gordon
theory, then generalize to higher correlation functions in the free theory.
Finally, we will consider $\phi^4$ theory, in which we can perform a
perturbation expansion to obtain the same Feynman rules as in Section~4.4.

Consider first a noninteracting real-valued scalar field:
\begin{equation}
  S_0 = \int d^4x\ \mathcal{L}_0
  = \int d^4x\ \Bigl[ \tfrac{1}{2}(\partial_\mu\phi)^2
    - \tfrac{1}{2}m^2\phi^2 \Bigr].
  \label{eq:9.19}
\end{equation}
Since $\mathcal{L}_0$ is quadratic in $\phi$, the functional integrals in
\eqref{eq:9.18} take the form of generalized, infinite-dimensional Gaussian
integrals.  We will therefore be able to evaluate the functional integrals
exactly.

Since this is our first functional integral computation, we will do it in a
very explicit, but ugly, way.  We must first define the integral
$\int\mathcal{D}\phi$ over field configurations.  To do this, we use the method
of Eq.~\eqref{eq:9.4} in considering the continuous integral as a limit of a
large but finite number of integrals.  We thus replace the variables
$\phi(x)$ defined on a continuum of points by variables $\phi(x_i)$ defined at
the points $x_i$ of a square lattice.  Let the lattice spacing be $\epsilon$,
let the four-dimensional spacetime volume be $L^4$, and define
\begin{equation}
  \mathcal{D}\phi = \prod_i d\phi(x_i),
  \label{eq:9.20}
\end{equation}
up to an irrelevant overall constant.

The field values $\phi(x_i)$ can be represented by a discrete Fourier series:
\begin{equation}
  \phi(x) = \sum_{k_n} \frac{1}{\sqrt{V}}\ \phi(k_n)\,e^{ik_n\cdot x},
  \qquad \phi(k_n) = \frac{1}{\sqrt{V}}\int d^4x\ \phi(x)\,e^{-ik_n\cdot x},
\end{equation}
where $k_n^\mu = 2\pi n^\mu/L$, with $n^\mu$ an integer, $|k_n^\mu| < \pi/\epsilon$,
and $V = L^4$.  The Fourier coefficients $\phi(k)$ are complex.  However,
$\phi(x)$ is real, and so these coefficients must obey the constraint
$\phi^{*}(k) = \phi(-k)$.  We will consider the real and imaginary parts of the
$\phi(k_n)$ with $k_n^0 > 0$ as independent variables.  The change of variables
from the $\phi(x_i)$ to these new variables $\phi(k_n)$ is a unitary
transformation, so we can rewrite the integral as
\begin{equation}
  \mathcal{D}\phi(x) = \Bigl[ \prod_{k_n^0>0} d\,\mathrm{Re}\,\phi(k_n)\,
    d\,\mathrm{Im}\,\phi(k_n) \Bigr] .
\end{equation}
Later, we will take the limit $L\to\infty$, $\epsilon\to 0$.  The effect of
this limit is to convert discrete, finite sums over $k_n$ to continuous
integrals over $k$:
\begin{equation}
  \frac{1}{V}\sum_{k_n} \to \int \frac{d^4k}{(2\pi)^4} .
  \label{eq:9.22}
\end{equation}
In the following discussion, this limit will produce Feynman perturbation
theory in the form derived in Part~I.  We will not eliminate the infrared and
ultraviolet divergences of Feynman diagrams that we encountered in Chapter~6,
but at least the functional integral introduces no new types of singular
behavior.

Now, in the free theory, the exponent of \eqref{eq:9.18} is
\begin{equation}
  S_0 = \int \frac{d^4k}{(2\pi)^4}\ \Bigl[ \phi^{*}(k)
    \bigl( \tfrac{1}{2}k^2 - \tfrac{1}{2}m^2 \bigr)\phi(k) \Bigr] .
\end{equation}
The Gaussian integral over $\phi$ can be performed explicitly.  For a
quadratic form, the result is the inverse of the quadratic coefficient.  We
find, for the two-point function,
\begin{equation}
  \langle 0|T\phi(x_1)\phi(x_2)|0\rangle = D_F(x_1-x_2)
  = \int \frac{d^4k}{(2\pi)^4}\ \frac{i}{k^2-m^2+i\epsilon}\
    e^{-ik\cdot(x_1-x_2)} ,
  \label{eq:9.28}
\end{equation}
which is the familiar Feynman propagator of the Klein--Gordon theory.  The
$i\epsilon$ term arises from the convergence factor of the functional integral,
as we noted below Eq.~\eqref{eq:9.5}.  For the four-point function, we
evaluate the Gaussian integral and obtain
\begin{equation}
  \langle 0|T\phi(x_1)\phi(x_2)\phi(x_3)\phi(x_4)|0\rangle
  = D_F(x_1-x_2)D_F(x_3-x_4)
    + D_F(x_1-x_3)D_F(x_2-x_4)
    + D_F(x_1-x_4)D_F(x_2-x_3) ,
  \label{eq:9.29}
\end{equation}
in agreement with Wick's theorem.

We are now ready to move from the free Klein--Gordon theory to $\phi^4$
theory.  Add to $\mathcal{L}_0$ a $\phi^4$ interaction:
\begin{equation}
  \mathcal{L} = \mathcal{L}_0 - \frac{\lambda}{4!}\phi^4 .
\end{equation}
Assuming that $\lambda$ is small, we can expand
\begin{equation}
  \exp\Bigl[ i\int d^4x\ \mathcal{L} \Bigr] = \exp\Bigl[
    i\int d^4x\ \mathcal{L}_0 \Bigr]\ \Bigl[ 1 - \frac{i\lambda}{4!}
      \int d^4x\ \phi^4 + \cdots \Bigr] .
\end{equation}
Making this expansion in both the numerator and the denominator of
\eqref{eq:9.18}, we see that each is (aside from the constant factor
Eq.~(9.23), which again cancels) expressed entirely in terms of free-field
correlation functions.  Moreover, since
$i\int d^4x\ \mathcal{L}_{\mathrm{int}} = -i\lambda/4!\int d^4x\ \phi^4$, we
obtain exactly the same expansion as in Eq.~\eqref{eq:4.31}.  We can express
both the numerator and the denominator in terms of Feynman diagrams, with the
fundamental interaction again given by the vertex
\begin{equation}
  = -i\lambda\,(2\pi)^4\delta^{(4)}(\textstyle\sum k_i) .
  \label{eq:9.30}
\end{equation}
All of the combinatorics work the same as in Section~4.4.  In particular, the
disconnected vacuum bubble diagrams exponentiate and factor from the
numerator of \eqref{eq:9.18}, and are canceled by the denominator, just as in
Eq.~\eqref{eq:4.31}.

The vertex rule for $\phi^4$ theory follows from the Lagrangian in an
exceedingly simple way, and this simple procedure will turn out to be valid
for other quantum field theories as well.  Once the quadratic terms in the
Lagrangian are properly understood and the propagators of the theory are
computed, the vertices can be read directly from the Lagrangian as the
coefficients of the cubic and higher-order terms.

\subsection{Functional Derivatives and the Generating Functional}

To conclude this section, we will now introduce a slicker, more formal, method
for computing correlation functions.  This method, based on an object called
the \emph{generating functional}, avoids the awkward Fourier expansions of the
preceding derivation.

First we define the functional derivative, $\delta/\delta J(x)$, as follows.
The functional derivative obeys the basic axiom (in four dimensions)
\begin{equation}
  \frac{\delta}{\delta J(x)}\,J(y) = \delta^{(4)}(x-y)
  \qquad\text{or}\qquad
  \frac{\delta}{\delta J(x)}\,\int d^4y\ J(y)\phi(y) = \phi(x) .
  \label{eq:9.31}
\end{equation}
This definition is the natural generalization, to continuous functions, of the
rule for discrete vectors,
\begin{equation}
  \frac{\partial}{\partial x_i}\,x_j = \delta_{ij} .
\end{equation}
To take functional derivatives of more complicated functionals we simply use
the ordinary rules for derivatives of composite functions.  For example,
\begin{equation}
  \frac{\delta}{\delta J(x)}\ \exp\Bigl[ i\int d^4y\ J(y)\phi(y) \Bigr]
  = i\phi(x)\ \exp\Bigl[ i\int d^4y\ J(y)\phi(y) \Bigr] .
  \label{eq:9.32}
\end{equation}
When the functional depends on the derivative of $J$, we integrate by parts
before applying the functional derivative:
\begin{equation}
  \frac{\delta}{\delta J(x)}\ \int d^4y\ \partial_\mu J(y)\,\phi(y)
  = -\partial_\mu\phi(x) .
  \label{eq:9.33}
\end{equation}

The basic object of this formalism is the \emph{generating functional of
correlation functions}, $Z[J]$.  (Some authors call it $W[J]$.)  In a scalar
field theory, $Z[J]$ is defined as
\begin{equation}
  Z[J] = \int \mathcal{D}\phi\ \exp\Bigl[ i\int d^4x\ \bigl(
    \mathcal{L} + J(x)\phi(x) \bigr) \Bigr] .
  \label{eq:9.34}
\end{equation}
This is a functional integral over $\phi$ in which we have added to
$\mathcal{L}$ in the exponent a source term, $J(x)\phi(x)$.

Correlation functions of the Klein--Gordon field theory can be simply
computed by taking functional derivatives of the generating functional.  For
example, the two-point function is
\begin{equation}
  \langle 0|T\phi(x_1)\phi(x_2)|0\rangle
  = \frac{1}{Z_0}\ \Bigl( -i\frac{\delta}{\delta J(x_1)} \Bigr)\
    \Bigl( -i\frac{\delta}{\delta J(x_2)} \Bigr)\ Z[J]\Big|_{J=0},
  \label{eq:9.35}
\end{equation}
where $Z_0 = Z[J=0]$.  Each functional derivative brings down a factor of
$\phi$ in the numerator of $Z[J]$; setting $J=0$, we recover expression
\eqref{eq:9.18}.  To compute higher correlation functions we simply take more
functional derivatives.

Formula \eqref{eq:9.35} is useful because, in a free field theory, $Z[J]$ can
be rewritten in a very explicit form.  Consider the exponent of \eqref{eq:9.34}
in the free Klein--Gordon theory.  Integrating by parts, we obtain
\begin{equation}
  \int d^4x\ \bigl[ \mathcal{L}_0(\phi) + J\phi \bigr]
  = \int d^4x\ \Bigl[ \tfrac{1}{2}\phi(-\partial^2-m^2+i\epsilon)\phi
    + J\phi \Bigr].
  \label{eq:9.36}
\end{equation}
(The $i\epsilon$ is a convergence factor for the functional integral, as we
discussed below Eq.~(9.23).)  We can complete the square by
introducing a shifted field,
\begin{equation}
  \phi'(x) = \phi(x) - i\int d^4y\ D_F(x-y)\,J(y) .
\end{equation}
Making this substitution and using the fact that $D_F$ is a Green's function
of the Klein--Gordon operator, we find that \eqref{eq:9.36} becomes
\begin{equation}
  \int d^4x\ \bigl[ \mathcal{L}_0(\phi)+J\phi \bigr]
  = \int d^4x\ \Bigl[ \tfrac{1}{2}\phi'(-\partial^2-m^2+i\epsilon)\phi'
    - \tfrac{1}{2}J(-iD_F)J \Bigr] .
\end{equation}
More symbolically, we could write the change of variables as
\begin{equation}
  \phi' = \phi + (-\partial^2-m^2+i\epsilon)^{-1}J,
  \label{eq:9.37}
\end{equation}
and the result
\begin{equation}
  \int d^4x\ \bigl[ \mathcal{L}_0(\phi)+J\phi \bigr]
  = \int d^4x\ \bigl[ \tfrac{1}{2}\phi'(-\partial^2-m^2+i\epsilon)\phi'
    - \tfrac{1}{2}J(-\partial^2-m^2+i\epsilon)^{-1}J \bigr].
  \label{eq:9.38}
\end{equation}
Now change variables from $\phi$ to $\phi'$ in the functional integral of
\eqref{eq:9.34}.  This is just a shift, and so the Jacobian of the
transformation is 1.  The result is
\begin{equation}
  Z[J] = \int \mathcal{D}\phi'\ \exp\Bigl[ i\int d^4x\ \mathcal{L}_0(\phi')
    \Bigr]\ \exp\Bigl[ -\frac{i}{2}\int d^4x\,d^4y\ J(x)
      \bigl[ -iD_F(x-y) \bigr]J(y) \Bigr] .
\end{equation}
The second exponential factor is independent of $\phi'$, while the remaining
integral over $\phi'$ is precisely $Z_0$.  Thus the generating functional of
the free Klein--Gordon theory is simply
\begin{equation}
  Z[J] = Z_0\ \exp\Bigl[ \frac{i}{2}\int d^4x\,d^4y\ J(x)\,
    \frac{iD_F(x-y)}{i}\,J(y) \Bigr]
  = Z_0\ \exp\Bigl[ -\frac{1}{2}\int d^4x\,d^4y\ J(x)\,D_F(x-y)\,J(y)
    \Bigr].
  \label{eq:9.39}
\end{equation}
Let us use Eqs.~\eqref{eq:9.39} and \eqref{eq:9.35} to compute some
correlation functions.  The two-point function is
\begin{equation}
  \begin{aligned}
    \langle 0|T\phi(x_1)\phi(x_2)|0\rangle
    &= \frac{1}{Z_0}\ \frac{\delta}{\delta J(x_1)}\ \frac{\delta}
      {\delta J(x_2)}\ \exp\Bigl[ -\frac{1}{2}\int d^4x\,d^4y\
        J(x)D_F(x-y)J(y) \Bigr]\Big|_{J=0}\\
    &= D_F(x_1-x_2) .
  \end{aligned}
  \label{eq:9.40}
\end{equation}
Taking one derivative brings down two identical terms; the second derivative
gives several terms, but only when it acts on the outside factor do we get a
term that survives when we set $J=0$.

It is instructive to work out the four-point function by this method as well.
In order to fit the computation in a reasonable amount of space, let us
abbreviate arguments of functions as subscripts: $\phi_i = \phi(x_i)$,
$J_i = J(x_i)$, $D_{ij} = D_F(x_i-x_j)$, and so on.  Repeated subscripts will
be integrated over implicitly.  The four-point function is then
\begin{equation}
  \begin{aligned}
    \langle 0|T\phi_1\phi_2\phi_3\phi_4|0\rangle
    &= \frac{\delta}{\delta J_1}\frac{\delta}{\delta J_2}
      \frac{\delta}{\delta J_3}\frac{\delta}{\delta J_4}\
      \exp\Bigl[ -\frac{1}{2}J_iD_{ij}J_j \Bigr]\Big|_{J=0}\\
    &= D_{12}D_{34} + D_{13}D_{24} + D_{14}D_{23} ,
  \end{aligned}
  \label{eq:9.41}
\end{equation}
in agreement with \eqref{eq:9.29}.  The rules for differentiating the
exponential give rise to the same familiar pattern: We get one term for each
possible way of contracting the four points in pairs, with a factor of $D_F$
for each contraction.

The generating functional method used just above to construct the
correlations of a free field theory can be used as well to represent the
correlation functions of an interacting field theory.  Formula \eqref{eq:9.35}
is independent of whether the theory is free or interacting.  The factor
$Z[J=0]$ is nontrivial in the case of an interacting field theory, but it
simply gives the denominator of Eq.~\eqref{eq:9.18}, that is, the sum of
vacuum diagrams.  Again from this approach, the combinatoric issues in the
evaluation of correlation functions are the same as in Section~4.4.

\section{The Analogy Between Quantum Field Theory and Statistical Mechanics}

Let us now pause from the technical aspects of this discussion to consider some
implications of the formulae we have derived.  To begin, let us summarize the
formal conclusions of the previous section in the following way: For a field
theory governed by the Lagrangian $\mathcal{L}$, the generating functional of
correlation functions is
\begin{equation}
  Z[J] = \int \mathcal{D}\phi\ \exp\Bigl[ i\int d^4x\
    \bigl( \mathcal{L} + J\phi \bigr) \Bigr] .
  \label{eq:9.42}
\end{equation}
The time variable of integration in the exponent runs from $-T$ to $T$, with
$T\to\infty(1-i\epsilon)$.  A correlation function such as \eqref{eq:9.18} is
reproduced by writing
\begin{equation}
  \langle \Omega|T\phi(x_1)\phi(x_2)|\Omega\rangle
  = Z[J]^{-1}\ \Bigl( -i\frac{\delta}{\delta J(x_1)} \Bigr)\
    \Bigl( -i\frac{\delta}{\delta J(x_2)} \Bigr)\ Z[J]\Big|_{J=0} .
  \label{eq:9.43}
\end{equation}
The generating functional \eqref{eq:9.42} is reminiscent of the partition
function of statistical mechanics.  It has the same general structure of an
integral over all possible configurations of an exponential statistical weight.
The source $J(x)$ plays the role of an external field.  In fact, our method of
computing correlation functions by differentiating with respect to $J(x)$
mimics the trick often used in statistical mechanics of computing correlation
functions by differentiating with respect to such variables as the pressure or
the magnetic field.

This analogy can be made more precise by manipulating the time variable of
integration in \eqref{eq:9.42}.  The derivation of the functional integral
formula implied that the time integration was slightly tipped into the complex
plane, in just the direction to permit the contour to be rotated clockwise
onto the imaginary axis.  We have already noted (below Eq.~(9.23)) that
the original infinitesimal rotation gives the correct $i\epsilon$ prescription
to produce the Feynman propagator.  The finite rotation is the analogue in
configuration space of the Wick rotation of the time component of momentum
illustrated in Fig.~\ref{fig:6.1}.  Like the Wick rotation in a momentum
integral, this Wick rotation of the time coordinate $t\to -ix^0$ produces a
Euclidean 4-vector product:
\begin{equation}
  x^2 = t^2 - |\mathbf{x}|^2 \to -(x^0)^2 - |\mathbf{x}|^2 = -|x_E|^2 .
  \label{eq:9.44}
\end{equation}
It is possible to show, by manipulating the expression for each Feynman
diagram, that the analytic continuation of the time variables in any Green's
function of a quantum field theory produces a correlation function invariant
under the rotational symmetry of four-dimensional Euclidean space.  This Wick
rotation inside the functional integral demonstrates this same conclusion in a
more general way.

To understand what we have achieved by this rotation, consider the example of
$\phi^4$ theory.  The action of $\phi^4$ theory coupled to sources is
\begin{equation}
  \int d^4x\ \bigl( \mathcal{L} + J\phi \bigr) = \int d^4x\ \Bigl[
    \tfrac{1}{2}(\partial_\mu\phi)^2 - \tfrac{1}{2}m^2\phi^2
      - \frac{\lambda}{4!}\phi^4 + J\phi \Bigr] .
\end{equation}
After the Wick rotation $t\to -ix^0$, this becomes
\begin{equation}
  i\int d^4x\ \bigl( \mathcal{L}+J\phi \bigr) \to -\int d^4x_E\ \Bigl[
    \tfrac{1}{2}(\nabla_E\phi)^2 + \tfrac{1}{2}m^2\phi^2
      + \frac{\lambda}{4!}\phi^4 - J\phi \Bigr] ,
\end{equation}
which is precisely the exponent appearing in the statistical partition
function of the magnetic system discussed in Chapter~8.  The identification is
completed by noting that the Euclidean field theory in $d$ spatial dimensions
corresponds to the statistical system at a temperature given by the inverse
period of the Euclidean time coordinate, which is taken to infinity.  This is
the precise sense in which a quantum field theory in $d+1$ spacetime
dimensions corresponds to a statistical system in $d$ spatial dimensions.  The
analogy is fully developed in Problem~9.2, where we evaluate the finite-
temperature (periodic Euclidean time) version of the path integral.

\section{Quantization of the Electromagnetic Field}

We now turn to the quantization of the electromagnetic field, which we have
been avoiding since Section~4.8.  In this section we will apply the functional
integral method to the quantization of the photon.  The new difficulty, of
course, is gauge invariance.

Let us begin by attempting to write down the generating functional for the
electromagnetic field in the same way as for a scalar field:
\begin{equation}
  Z[J] = \int \mathcal{D}A\ \exp\Bigl[ i\int d^4x\ \bigl(
    \mathcal{L} + J_\mu A^\mu \bigr) \Bigr] ,
  \label{eq:9.47}
\end{equation}
where
\begin{equation}
  \mathcal{L} = -\frac{1}{4}(F_{\mu\nu})^2
\end{equation}
is the free electromagnetic Lagrangian.  If we work in the free-field theory,
we can perform the functional integral over $A$ exactly, just as we did for the
scalar field.  The quadratic form in the exponent is
\begin{equation}
  \int d^4x\ \Bigl[ -\frac{1}{4}(F_{\mu\nu})^2 + J_\mu A^\mu \Bigr]
  = \int d^4x\ \Bigl[ \tfrac{1}{2}A^\mu\bigl( g_{\mu\nu}\partial^2
    - \partial_\mu\partial_\nu \bigr)A^\nu + J_\mu A^\mu \Bigr] .
  \label{eq:9.52}
\end{equation}
The operator $(g_{\mu\nu}\partial^2 - \partial_\mu\partial_\nu)$ has no
inverse: it annihilates any function proportional to $\partial_\mu\Lambda(x)$,
since
\begin{equation}
  \bigl( g_{\mu\nu}\partial^2 - \partial_\mu\partial_\nu \bigr)
    \partial^\nu\Lambda = 0 .
\end{equation}
This is, of course, just the gauge invariance of the photon: the action is
unchanged by the transformation $A_\mu\to A_\mu + \partial_\mu\Lambda$.  The
quadratic form is singular, and the functional integral over $A$ is not
well-defined as it stands.  Physically, we must not integrate over the gauge
degrees of freedom; we should integrate over only the physical, gauge-inequivalent
field configurations.

The solution to this problem is the \emph{Faddeev--Popov} procedure.  We
insert a factor of 1 into the functional integral, in the form of an integral
over the gauge transformation group.  Consider the functional integral
\begin{equation}
  \int \mathcal{D}A\ e^{iS[A]} ,
\end{equation}
where $S[A]$ is gauge invariant.  We insert
\begin{equation}
  1 = \Delta[A]\ \int \mathcal{D}\alpha\ \delta\bigl( G(A^\alpha) \bigr) ,
  \label{eq:9.53}
\end{equation}
where $G(A)$ is a gauge-fixing function (which vanishes for exactly one gauge
orbit representative), $A^\alpha$ is the gauge-transformed field, and
$\Delta[A]$ is the Faddeev--Popov determinant, defined so that the integral is
1.  Using the gauge invariance of the action and of the measure, we can
evaluate the functional integral:
\begin{equation}
  \int \mathcal{D}A\ e^{iS[A]} = \det\Bigl( \frac{\delta G}{\delta\alpha}
    \Bigr)\ \Bigl( \int \mathcal{D}\alpha \Bigr)\
    \int \mathcal{D}A\ e^{iS[A]}\ \delta\bigl( G(A) \bigr) .
  \label{eq:9.54}
\end{equation}
The functional integral over $A$ is now restricted by the delta function to
physically inequivalent field configurations, as desired.  The divergent
integral over $\alpha(x)$ simply gives an infinite multiplicative factor.

To go further we must specify a gauge-fixing function $G(A)$.  We choose the
general class of functions
\begin{equation}
  G(A) = \partial^\mu A_\mu(x) - \omega(x) ,
  \label{eq:9.55}
\end{equation}
where $\omega(x)$ can be any scalar function.  Setting this $G(A)$ equal to
zero gives a generalization of the Lorentz gauge condition.  The functional
determinant is the same as in Lorentz gauge,
$\det(\delta G(A^\alpha)/\delta\alpha) = \det(\partial^2/e)$.  Thus the
functional integral becomes
\begin{equation}
  \int \mathcal{D}A\ e^{iS[A]} = \det\Bigl( \frac{1}{e}\partial^2 \Bigr)\
    \Bigl( \int \mathcal{D}\alpha \Bigr)\ \int \mathcal{D}A\ e^{iS[A]}\
    \delta\bigl( \partial^\mu A_\mu - \omega(x) \bigr) .
\end{equation}
This equality holds for any $\omega(x)$, so it will also hold if we replace the
right-hand side with any properly normalized linear combination involving
different functions $\omega(x)$.  For our final trick, we will integrate over
all $\omega(x)$, with a Gaussian weighting function centered on $\omega=0$.
The above expression is thus equal to
\begin{equation}
  \begin{aligned}
    \int \mathcal{D}A\ e^{iS[A]} &= N(\xi)\ \det\Bigl( \frac{1}{e}\partial^2
      \Bigr)\ \Bigl( \int \mathcal{D}\alpha \Bigr)\ \int \mathcal{D}A\
      e^{iS[A]}\ \exp\Bigl[ -\frac{i}{2\xi}\int d^4x\
        (\partial^\mu A_\mu)^2 \Bigr] ,
  \end{aligned}
  \label{eq:9.56}
\end{equation}
where $N(\xi)$ is an unimportant normalization constant and we have used the
delta function to perform the integral over $\omega$.  We can choose $\xi$ to
be any finite constant.  Effectively, we have added a new term
$-(\partial^\mu A_\mu)^2/2\xi$ to the Lagrangian.

So far we have worked only with the denominator of our formula for correlation
functions,
\begin{equation}
  \langle\Omega|T\mathcal{O}(A)|\Omega\rangle
  = \lim_{T\to\infty(1-i\epsilon)}\
    \frac{\int\mathcal{D}A\ \mathcal{O}(A)\ \exp\bigl[
      i\int_{-T}^{T}d^4x\ \mathcal{L} \bigr]}
      {\int\mathcal{D}A\ \exp\bigl[ i\int_{-T}^{T}d^4x\ \mathcal{L} \bigr]} .
\end{equation}
The same manipulations can also be performed on the numerator, provided that
the operator $\mathcal{O}(A)$ is gauge invariant.  (If it is not, the variable
change from $A$ to $A^\alpha$ preceding Eq.~\eqref{eq:9.54} does not work.)
Assuming that $\mathcal{O}(A)$ is gauge invariant, we find for its correlation
function
\begin{equation}
  \langle\Omega|T\mathcal{O}(A)|\Omega\rangle
  = \lim_{T\to\infty(1-i\epsilon)}\
    \frac{\int\mathcal{D}A\ \mathcal{O}(A)\ \exp\bigl[ i\int_{-T}^{T}d^4x\
      \bigl( \mathcal{L} - \frac{1}{2\xi}(\partial^\mu A_\mu)^2 \bigr)
      \bigr]}
      {\int\mathcal{D}A\ \exp\bigl[ i\int_{-T}^{T}d^4x\
        \bigl( \mathcal{L} - \frac{1}{2\xi}(\partial^\mu A_\mu)^2 \bigr)
        \bigr]} .
  \label{eq:9.57}
\end{equation}
The awkward constant factors in \eqref{eq:9.56} have canceled; the only trace
left by this whole process is the extra $\xi$-term that is added to the
action.

At the beginning of this section, in Eq.~\eqref{eq:9.52}, we saw that we
could not obtain a sensible photon propagator from the action $S[A]$.  With
the new $\xi$-term, however, that equation becomes
\begin{equation}
  \bigl( -k^2 g_{\mu\nu} + \bigl( 1 - \tfrac{1}{\xi} \bigr)k_\mu k_\nu \bigr)
    \tilde{D}^{\nu\rho}(k) = -i\delta^\rho_\mu ,
\end{equation}
which has the solution
\begin{equation}
  \tilde{D}^{\mu\nu}(k) = \frac{-i}{k^2+i\epsilon}\ \Bigl[
    g^{\mu\nu} - (1-\xi)\frac{k^\mu k^\nu}{k^2} \Bigr] .
  \label{eq:9.58}
\end{equation}
This is our desired expression for the photon propagator.  The $i\epsilon$
term in the denominator arises exactly as in the Klein--Gordon case.  Note the
overall minus sign relative to the Klein--Gordon propagator, which was already
evident in Eq.~\eqref{eq:9.52}.

In practice one usually chooses a specific value of $\xi$ when making
computations.  Two choices that are often convenient are
\begin{equation}
  \xi = 0 \qquad \text{Landau gauge};\qquad
  \xi = 1 \qquad \text{Feynman gauge} .
\end{equation}
So far in this book we have always used Feynman gauge.%
\footnote{Other choices of $\xi$ may be useful in specific applications; for
example, in certain problems of bound states in QED, the Yennie gauge,
$\xi = 3$, produces a cancellation that is otherwise difficult to make
explicit.  See H.~M.~Fried and D.~R.~Yennie, \emph{Phys.\ Rev.} \textbf{112},
1391 (1958).}

The Faddeev--Popov procedure guarantees that the value of any correlation
function of gauge-invariant operators computed from Feynman diagrams will be
independent of the value of $\xi$ used in the calculation (as long as the same
value of $\xi$ is used consistently).  In the case of QED, it is not difficult
to prove this $\xi$-independence directly.  Notice in Eq.~\eqref{eq:9.58}
that $\xi$ multiplies a term in the photon propagator proportional to
$k^\mu k^\nu$.  According to the Ward--Takahashi identity \eqref{eq:7.68}, the
replacement in a Green's function of any photon propagator by $k^\mu k^\nu$
yields zero, except for terms involving external off-shell fermions.  These
terms are equal and opposite for particle and antiparticle and vanish when the
fermions are grouped into gauge-invariant combinations.

To complete our treatment of the quantization of the electromagnetic field, we
need one additional ingredient.  In Chapters~5 and~6, we computed $S$-matrix
elements for QED from the correlation functions of non-gauge-invariant
operators $\psi(x)$, $\bar{\psi}(x)$, and $A^\mu(x)$.  We will now argue that
the $S$-matrix elements are given correctly by this procedure.  Since the
$S$-matrix is defined between asymptotic states, we can compute $S$-matrix
elements in a formalism in which the coupling constant is turned off
adiabatically in the far past and far future.  In the zero coupling limit,
there is a clean separation between gauge-invariant and gauge-variant states.
Single-particle states containing one electron, one positron, or one
transversely polarized photon are gauge-invariant, while states with timelike
and longitudinal photon polarizations transform under gauge motions.  We can
thus define a gauge-invariant $S$-matrix in the following way: Let $S_{FP}$ be
the $S$-matrix between general asymptotic states, computed from the
Faddeev--Popov procedure.  This matrix is unitary but not gauge-invariant.
Let $P_0$ be a projection onto the subspace of the space of asymptotic states
in which all particles are either electrons, positrons, or transverse photons.
Then let
\begin{equation}
  S = P_0\,S_{FP}\,P_0 .
  \label{eq:9.59}
\end{equation}
This $S$-matrix is gauge invariant by construction, because it is projected
onto gauge-invariant states.  It is not obvious that it is unitary.  However,
we addressed this issue in Section~5.5.  We showed there that any matrix
element $\mathcal{M}^\mu\epsilon^{*}_\mu$ for photon emission satisfies
\begin{equation}
  \sum_{\epsilon=1,2} \bigl| \mathcal{M}^\mu\epsilon^{*}_\mu \bigr|^2
  = -g_{\mu\nu}\mathcal{M}^\mu\mathcal{M}^{\nu*} ,
  \label{eq:9.60}
\end{equation}
where the sum on the left-hand side runs only over transverse polarizations.
The same argument applies if $\mathcal{M}^\mu$ and $\mathcal{M}^{\nu*}$ are
distinct amplitudes, as long as they satisfy the Ward identity.  This is
exactly the information we need to see that
\begin{equation}
  SS^{\dagger} = P_0\,S_{FP}\,P_0\,S_{FP}^{\dagger}\,P_0
    = P_0\,S_{FP}\,S_{FP}^{\dagger}\,P_0 .
  \label{eq:9.61}
\end{equation}
Now we can use the unitarity of $S_{FP}$ to see that $S$ is unitary,
$S^{\dagger}S = 1$, on the subspace of gauge-invariant states.  It is easy to
check explicitly that the formula \eqref{eq:9.59} for the $S$-matrix is
independent of $\xi$: The Ward identity implies that any QED matrix element
with all external fermions on-shell is unchanged if we add to the photon
propagator $D^{\mu\nu}(q)$ any term proportional to $q^\mu q^\nu$.

\section{Functional Quantization of Spinor Fields}

The functional methods that we have used so far allow us to compute, using
Eq.~\eqref{eq:9.18} or \eqref{eq:9.35}, correlation functions involving fields
that obey canonical commutation relations.  To generalize these methods to
include spinor fields, which obey canonical anticommutation relations, we must
do something different: We must represent even the classical fields by
\emph{anticommuting numbers}.

\subsection{Anticommuting Numbers}

We will define anticommuting numbers (also called Grassmann numbers) by giving
algebraic rules for manipulating them.  These rules are formal and might seem
ad hoc.  We will justify them by showing that they lead to the familiar
quantum theory of the Dirac equation.

The basic feature of anticommuting numbers is that they anticommute.  For any
two such numbers $\theta$ and $\eta$,
\begin{equation}
  \theta\eta = -\eta\theta .
  \label{eq:9.62}
\end{equation}
In particular, the square of any Grassmann number is zero:
\begin{equation}
  \theta^2 = 0 .
\end{equation}
(This fact makes algebra extremely easy.)  A product $(\theta\eta)$ of two
Grassmann numbers commutes with other Grassmann numbers.  We will also wish to
add Grassmann numbers, and to multiply them by ordinary numbers; these
operations have all the properties of addition and scalar multiplication in
any vector space.

The main thing we want to do with anticommuting numbers is integrate over
them.  To define functional integration, we do not need general definite
integrals of these parameters, but only the analog of $\int dx$.  So let us
define the integral of a general function $f$ of a Grassmann variable
$\theta$, over the complete range of $\theta$:
\begin{equation}
  \int d\theta\ f(\theta) = \int d\theta\ (A + B\theta) .
\end{equation}
In general, $f(\theta)$ can be expanded in a Taylor series, which terminates
after two terms since $\theta^2 = 0$.  The integral should be linear in $f$;
thus it must be a linear function of $A$ and $B$.  Its value is fixed by one
additional property: In our analysis of bosonic functional integrals (for
instance, in \eqref{eq:9.38} and \eqref{eq:9.54}), we made strong use of the
invariance of the integral to shifts of the integration variable.  We will see
in Section~9.6 that this shift invariance of the functional integral plays a
central role in the derivation of the quantum mechanical equations of motion
and conservation laws, and thus must be considered a fundamental aspect of the
formalism.  We must, then, demand this same property for integrals over
$\theta$.  Invariance under the shift $\theta\to\theta+\eta$ yields the
condition
\begin{equation}
  \int d\theta\ (A+B\theta) = \int d\theta\ \bigl( (A+B\eta) + B\theta \bigr) .
\end{equation}
The shift changes the constant term, but leaves the linear term unchanged.
The only linear function of $A$ and $B$ that has this property is a constant
(conventionally taken to be 1) times $B$, so we define
\begin{equation}
  \int d\theta\ (A + B\theta) = B .
  \label{eq:9.63}
\end{equation}
When we perform a multiple integral over more than one Grassmann variable, an
ambiguity in sign arises; we adopt the convention
\begin{equation}
  \int d\eta\,d\theta\ \theta\eta = +1 ,
  \label{eq:9.64}
\end{equation}
performing the innermost integral first.

Since the Dirac field is complex-valued, we will work primarily with complex
Grassmann numbers, which can be built out of real and imaginary parts in the
usual way.  It is convenient to define complex conjugation to reverse the
order of products, just like Hermitian conjugation of operators:
\begin{equation}
  (\theta\eta)^{*} = \eta^{*}\theta^{*} = -\theta^{*}\eta^{*} .
  \label{eq:9.65}
\end{equation}
To integrate over complex Grassmann numbers, let us define
\begin{equation}
  \theta = \frac{\theta_1 + i\theta_2}{\sqrt{2}},\qquad
  \theta^{*} = \frac{\theta_1 - i\theta_2}{\sqrt{2}} .
\end{equation}
We can now treat $\theta$ and $\theta^{*}$ as independent Grassmann numbers,
and adopt the convention $\int d\theta^{*}d\theta\ (\theta\theta^{*}) = 1$.

Let us evaluate a Gaussian integral over a complex Grassmann variable:
\begin{equation}
  \int d\theta^{*}d\theta\ e^{-\theta^{*}b\theta}
  = \int d\theta^{*}d\theta\ (1 - \theta^{*}b\theta)
  = \int d\theta^{*}d\theta\ (1 + \theta\theta^{*}b) = b .
  \label{eq:9.66}
\end{equation}
If $\theta$ were an ordinary complex number, this integral would equal
$2\pi/b$.  The factor of $2\pi$ is unimportant; the main difference with
anticommuting numbers is that the $b$ comes out in the numerator rather than
the denominator.  However, if there is an additional factor of $\theta\theta^{*}$
in the integrand, we obtain
\begin{equation}
  \int d\theta^{*}d\theta\ \theta\theta^{*}\ e^{-\theta^{*}b\theta}
  = 1 = 1/b .
  \label{eq:9.67}
\end{equation}
The extra $\theta\theta^{*}$ introduces a factor of $(1/b)$, just as it does in
an ordinary Gaussian integral.

To perform general Gaussian integrals in higher dimensions, we must first
prove that an integral over complex Grassmann variables is invariant under
unitary transformations.  Consider a set of $n$ complex Grassmann variables
$\theta_i$, and a unitary matrix $U$.  If $\theta'_i = U_{ij}\theta_j$, then
the transformation of the integration measure introduces the determinant of
$U$, which is a phase.  Together with the transformation of the integrand, the
phase cancels, and we find that the integral is invariant under unitary
transformations.

The general Gaussian integral over complex Grassmann variables is then
evaluated by diagonalizing the Hermitian matrix $M$ by a unitary
transformation, and using \eqref{eq:9.66} for each eigenvalue.  The result is
\begin{equation}
  \int \Bigl[ \prod_i d\theta^{*}_i d\theta_i \Bigr]\
    \exp\Bigl[ -\theta^{*}_i M_{ij}\theta_j \Bigr] = \det M .
  \label{eq:9.70}
\end{equation}
This is the fundamental formula for functional integrals over anticommuting
fields.  Note that the determinant appears in the numerator, in contrast to
the bosonic case, where the inverse determinant appears.

\subsection{The Dirac Propagator}

We are now ready to apply the functional method to the Dirac field.  The
generating functional for the free Dirac field is
\begin{equation}
  Z[\eta, \bar{\eta}] = \int \mathcal{D}\bar{\psi}\,\mathcal{D}\psi\
    \exp\Bigl[ i\int d^4x\ \Bigl( \bar{\psi}(i\slashed{\partial}-m)\psi
      + \bar{\eta}\psi + \bar{\psi}\eta \Bigr) \Bigr] ,
  \label{eq:9.71}
\end{equation}
where $\eta$ and $\bar{\eta}$ are Grassmann-valued sources.  Completing the
square as in the scalar case, and using the Gaussian integral
\eqref{eq:9.70}, we find
\begin{equation}
  Z[\eta, \bar{\eta}] = Z_0\ \exp\Bigl[ -\int d^4x\,d^4y\
    \bar{\eta}(x)\,S_F(x-y)\,\eta(y) \Bigr] ,
  \label{eq:9.74}
\end{equation}
where $S_F(x-y)$ is the Feynman propagator of the Dirac field,
\begin{equation}
  S_F(x-y) = \int \frac{d^4p}{(2\pi)^4}\ \frac{i(\slashed{p}+m)}
    {p^2-m^2+i\epsilon}\ e^{-ip\cdot(x-y)} .
\end{equation}
Correlation functions are obtained by differentiating with respect to the
sources.  For example,
\begin{equation}
  \langle 0|T\psi(x_1)\bar{\psi}(x_2)|0\rangle
  = \frac{1}{Z_0}\ \Bigl( -i\frac{\delta}{\delta\eta(x_1)} \Bigr)\
    \Bigl( -i\frac{\delta}{\delta\bar{\eta}(x_2)} \Bigr)\ Z[\eta,\bar{\eta}]
    \Big|_{\eta=\bar{\eta}=0}
  = S_F(x_1-x_2) .
  \label{eq:9.75}
\end{equation}
Plugging in formula \eqref{eq:9.74} for $Z[\eta,\bar{\eta}]$ and carefully
keeping track of the signs, we find that this expression is equal to the
Feynman propagator, $S_F(x_1-x_2)$.  Higher correlation functions can be
evaluated in a similar way.

\subsection{QED}

As we saw in Section~9.2 for the case of scalar fields, the functional
integral method allows us to read the Feynman rules for vertices directly from
the Lagrangian for an interacting field theory.  For the theory of Quantum
Electrodynamics, the full Lagrangian is
\begin{equation}
  \begin{aligned}
    \mathcal{L}_{\mathrm{QED}} &= \bar{\psi}(i\gamma^\mu\partial_\mu-m)\psi
      - \tfrac{1}{4}(F_{\mu\nu})^2\\
    &= \bar{\psi}(i\slashed{D}-m)\psi - \tfrac{1}{4}(F_{\mu\nu})^2
      - e\bar{\psi}\gamma^\mu\psi A_\mu\\
    &= \mathcal{L}_0 - e\bar{\psi}\gamma^\mu\psi A_\mu ,
  \end{aligned}
  \label{eq:9.76}
\end{equation}
where $D_\mu = \partial_\mu + ieA_\mu$ is the gauge-covariant derivative.

To evaluate correlation functions, we expand the exponential of the
interaction term:
\begin{equation}
  \exp\Bigl[ i\int \mathcal{L} \Bigr] = \exp\Bigl[ i\int \mathcal{L}_0 \Bigr]\
    \Bigl[ 1 - ie\int d^4x\ \bar{\psi}\gamma^\mu\psi A_\mu + \cdots \Bigr] .
\end{equation}
The two terms of the free Lagrangian yield the Dirac and electromagnetic
propagators derived in this section and the last:
\begin{equation}
  \psi\bar{\psi}: \quad \frac{i(\slashed{p}+m)}{p^2-m^2+i\epsilon};\qquad
  A^\mu A^\nu: \quad \frac{-ig^{\mu\nu}}{k^2+i\epsilon}
    \quad \text{(Feynman gauge)} .
\end{equation}
The interaction term gives the QED vertex,
\begin{equation}
  \text{vertex} = -ie\gamma^\mu .
\end{equation}
As in Chapter~4, we can rearrange these rules, performing the integrations
over vertex positions to obtain momentum-conserving delta functions, and using
these delta functions to perform most of the propagator momentum integrals.

The only remaining aspect of the QED Feynman rules is the placement of various
minus signs.  These signs are also built into the functional integral; for
example, interchanging $\theta_k$ and $\theta_j$ in Eq.~\eqref{eq:9.70} would
introduce a factor of $-1$.  We will see another example of a fermion minus
sign in the computation that follows.

\subsection{Functional Determinants}

Throughout this chapter we have encountered expressions that we wrote
formally as functional determinants.  To end this section, let us investigate
one of these objects more closely.  We will find that, at least in this case,
we can write the determinant explicitly as a sum of Feynman diagrams.

Consider the object
\begin{equation}
  \int \mathcal{D}\bar{\psi}\,\mathcal{D}\psi\ \exp\Bigl[ i\int d^4x\
    \bar{\psi}(i\slashed{D}-m)\psi \Bigr],
  \label{eq:9.77}
\end{equation}
where $D_\mu = \partial_\mu + ieA_\mu$ and $A_\mu(x)$ is a given external
background field.  Formally, this expression is a functional determinant:
\begin{equation}
  \begin{aligned}
    \int \mathcal{D}\bar{\psi}\,\mathcal{D}\psi\ \exp\Bigl[ i\int d^4x\
      \bar{\psi}(i\slashed{D}-m)\psi \Bigr]
    &= \det(i\slashed{D}-m) = \det(i\slashed{\partial}-m - e\slashed{A})\\
    &= \det(i\slashed{\partial}-m)\cdot\det\Bigl( 1 -
      \frac{1}{i\slashed{\partial}-m}(-e\slashed{A}) \Bigr) .
  \end{aligned}
\end{equation}
In the last form, the first term is an infinite constant.  The second term
contains the dependence of the determinant on the external field $A$.  We will
now show that this dependence is well defined and, in fact, is exactly
equivalent to the sum of vacuum diagrams.

To demonstrate this, we need only apply standard identities from linear
algebra.  First notice that, if a matrix $B$ has eigenvalues $b_i$, we can
write its determinant as
\begin{equation}
  \det B = \prod_i b_i = \exp\Bigl[ \sum_i \log b_i \Bigr]
  = \exp\Bigl[ \mathrm{Tr}\bigl( \log B \bigr) \Bigr],
  \label{eq:9.78}
\end{equation}
where the logarithm of a matrix is defined by its power series.  Applying this
identity to our determinant, and writing out the power series of the
logarithm, we obtain
\begin{equation}
  \det\Bigl( 1 - \frac{1}{i\slashed{\partial}-m}(-e\slashed{A}) \Bigr)
  = \exp\Bigl[ -\sum_{n=1}^{\infty} \frac{1}{n}\ \mathrm{Tr}\Bigl(
    \frac{1}{i\slashed{\partial}-m}(-e\slashed{A}) \Bigr)^n \Bigr] .
  \label{eq:9.79}
\end{equation}
Alternatively, we can evaluate this determinant by returning to expression
\eqref{eq:9.77} and using Feynman diagrams.  Expanding the interaction term,
we obtain the vertex rule $-ie\gamma^\mu$.  Our determinant is then equal to a
sum of Feynman diagrams,
\begin{equation}
  \text{(sum of all closed fermion-loop diagrams with } n \text{ external
    photon legs)} .
  \label{eq:9.80}
\end{equation}
The series exponentiates, since the disconnected diagrams are products of
connected pieces (with appropriate symmetry factors when a piece is
repeated).

Now let us evaluate the $n$th diagram in the exponent of \eqref{eq:9.79}.
There is a factor of $-1$ from the fermion loop, and a symmetry factor of $1/n$
since we could rotate the interactions around the diagram up to $n$ times
without changing it.  (The factor is not $1/n!$, because the cyclic order of
the interaction points is significant.)  The diagram is therefore
\begin{equation}
  -\int d^4x_1\cdots d^4x_n\ \mathrm{tr}\bigl( -ie\slashed{A}(x_1) \bigr)\
    S_F(x_2-x_1)\cdots\ \bigl( -ie\slashed{A}(x_n) \bigr)\ S_F(x_1-x_n)
  \label{eq:9.81}
\end{equation}
in exact agreement with \eqref{eq:9.79}, including the minus sign and the
symmetry factor.

The computation of functional determinants using Feynman diagrams is an
important tool, as we will see in Chapter~11.

\section{Symmetries in the Functional Formalism}

We have now seen that the quantum field theoretic correlation functions of
scalar, vector, and spinor fields can be computed from the functional integral,
completely bypassing the construction of the Hamiltonian, the Hilbert space of
states, and the equations of motion.  The functional integral formalism makes
the symmetries of the problem manifest; any invariance of the Lagrangian will
be an invariance of the quantum dynamics.%
\footnote{There are some subtle exceptions to this rule, which we will treat in
Chapter~19.}  However, we would like to be able to appeal also to the
conservation laws that follow from the quantum equations of motion, or to
these equations of motion themselves.  For example, the Ward identity, which
played a major role in our discussion of photons in QED (Section~5.5), is
essentially the conservation law of the electric charge current.  Since, as we
saw in Section~2.2, the conservation laws follow from symmetries of the
Lagrangian, one might guess that it is not difficult to derive these
conservation laws from the functional integral.  In this section we will see
how to do that.  We will see that the functional integral gives, in a most
direct way, a quantum generalization of Noether's theorem.  This result will
lead to the analogue of the Ward--Takahashi identity for any symmetry of a
general quantum field theory.

\subsection{Equations of Motion}

To prepare for this discussion, we should determine how the quantum equations
of motion follow from the functional integral formalism.  As a first problem
to study, let us examine the Green's functions of the free scalar field.  To
be specific, consider the three-point function:
\begin{equation}
  \langle\Omega|T\phi(x_1)\phi(x_2)\phi(x_3)|\Omega\rangle
  = Z^{-1}\int \mathcal{D}\phi\ e^{i\int d^4x\ \mathcal{L}}\
    \phi(x_1)\phi(x_2)\phi(x_3) ,
  \label{eq:9.82}
\end{equation}
where $\mathcal{L} = \tfrac{1}{2}(\partial_\mu\phi)^2 - \tfrac{1}{2}m^2\phi^2$
and $Z$ is a shorthand for $Z[J=0]$, the functional integral over the
exponential.  In classical mechanics, we would derive the equations of motion
by insisting that the action be stationary under an infinitesimal variation
\begin{equation}
  \phi(x) \to \phi'(x) = \phi(x) + \epsilon(x) .
  \label{eq:9.83}
\end{equation}
The appropriate generalization is to consider \eqref{eq:9.83} as an
infinitesimal change of variables.  A change of variables does not alter the
value of the integral.  Nor does a shift of the integration variable alter the
measure: $\mathcal{D}\phi' = \mathcal{D}\phi$.  Thus we can write
\begin{equation}
  \int \mathcal{D}\phi\ e^{i\int d^4x\ \mathcal{L}}\ \phi(x_1)\phi(x_2)
    \phi(x_3) = \int \mathcal{D}\phi'\ e^{i\int d^4x\ \mathcal{L}'}\
    \phi'(x_1)\phi'(x_2)\phi'(x_3) ,
\end{equation}
where $\phi' = \phi + \epsilon$.  Expanding this equation to first order in
$\epsilon$, we find
\begin{equation}
  \begin{aligned}
    0 &= \int \mathcal{D}\phi\ e^{i\int d^4x\ \mathcal{L}}\ \Bigl[
      i\int d^4x\ \epsilon(x)(-\partial^2-m^2)\phi(x)\ \phi(x_1)\phi(x_2)
        \phi(x_3)\\
    &\quad + \epsilon(x_1)\phi(x_2)\phi(x_3) + \phi(x_1)\epsilon(x_2)\phi(x_3)
      + \phi(x_1)\phi(x_2)\epsilon(x_3) \Bigr] .
  \end{aligned}
  \label{eq:9.84}
\end{equation}
The last three terms can be combined with the first by writing, for instance,
$\epsilon(x_1) = \int d^4x\ \epsilon(x)\,\delta(x-x_1)$.  Noting that the
right-hand side must vanish for any possible variation $\epsilon(x)$, we then
obtain
\begin{equation}
  \begin{aligned}
    0 &= \int \mathcal{D}\phi\ e^{i\int d^4x\ \mathcal{L}}\ \Bigl[
      i(\partial^2+m^2)\phi(x)\ \phi(x_1)\phi(x_2)\phi(x_3)\\
    &\quad + i\delta(x-x_1)\phi(x_2)\phi(x_3) + i\phi(x_1)\delta(x-x_2)
      \phi(x_3) + i\phi(x_1)\phi(x_2)\delta(x-x_3) \Bigr] .
  \end{aligned}
  \label{eq:9.85}
\end{equation}
A similar equation holds for any number of fields $\phi(x_i)$.

To see the implications of \eqref{eq:9.85}, let us specialize to the case of
one field $\phi(x_1)$ in \eqref{eq:9.82}.  Notice that the derivatives acting
on $\phi(x)$ can be pulled outside the functional integral.  Then, dividing
\eqref{eq:9.85} by $Z$ yields the identity
\begin{equation}
  (\partial^2 + m^2)\ \langle\Omega|T\phi(x)\phi(x_1)|\Omega\rangle
  = -i\delta(x-x_1) .
  \label{eq:9.86}
\end{equation}
The left-hand side of this relation is the Klein--Gordon operator acting on a
correlation function of $\phi(x)$.  The right-hand side is zero unless
$x = x_1$; that is, the correlation function satisfies the Klein--Gordon
equation except at the point where the arguments of the two $\phi$ fields
coincide.  The modification of the Klein--Gordon equation at this point is
called a \emph{contact term}.  In this simple case, the modification is hardly
unfamiliar to us; Eq.~\eqref{eq:9.86} merely says that the Feynman propagator
is a Green's function of the Klein--Gordon operator, as we originally showed
in Section~2.4.  We saw there that the delta function arises when the time
derivative in $\partial^2$ acts on the time-ordering symbol.  We will see below
that, quite generally in quantum field theory, the classical equations of
motion for fields are satisfied by all quantum correlation functions of those
fields, up to contact terms.

As an example, consider the identity that follows from \eqref{eq:9.85} for an
$(n+1)$-point correlation function of scalar fields:
\begin{equation}
  \begin{aligned}
    (\partial^2+m^2)\ \langle\Omega|T\phi(x)\phi(x_1)\cdots\phi(x_n)
      |\Omega\rangle
    &= \sum_{i=1}^{n}\ \langle\Omega|T\phi(x_1)\cdots\bigl(
      -i\delta(x-x_i) \bigr)\cdots\phi(x_n)|\Omega\rangle .
  \end{aligned}
  \label{eq:9.87}
\end{equation}
This identity says that the Klein--Gordon equation is obeyed by $\phi(x)$
inside any expectation value, up to contact terms associated with the time
ordering.  The result can also be derived from the Hamiltonian formalism using
the methods of Section~2.4, or, using the special properties of free-field
theory, by evaluating both sides of the equation using Wick's theorem.

As long as the functional measure is invariant under a shift of the
integration variable, we can repeat this argument and obtain the quantum
equations of motion for Green's functions for any theory of scalar, vector,
and spinor fields.  This is the reason why, in Eq.~\eqref{eq:9.63}, we took
the shift invariance to be the fundamental, defining property of the Grassmann
integral.

For a general field theory of a field $\phi(x)$, governed by the Lagrangian
$\mathcal{L}[\phi]$, the manipulations leading to \eqref{eq:9.84} give the
identity
\begin{equation}
  0 = \int \mathcal{D}\phi\ e^{i\int d^4x\ \mathcal{L}}\ \Bigl[
    i\frac{\delta S}{\delta\phi(x)}\ \phi(x_1)\cdots\phi(x_n)
    + \sum_{i=1}^{n} i\delta(x-x_i)\ \phi(x_1)\cdots\phi(x_n) \Bigr] ,
  \label{eq:9.88}
\end{equation}
and similar identities for correlation functions of $n$ fields.  By the rule
for functional differentiation \eqref{eq:9.31}, the derivative of the action
is
\begin{equation}
  \frac{\delta S}{\delta\phi(x)} = \frac{\partial\mathcal{L}}
    {\partial\phi(x)} - \partial_\mu
      \frac{\partial\mathcal{L}}{\partial(\partial_\mu\phi)} ;
\end{equation}
this is the quantity that equals zero by the Euler--Lagrange equation of
motion \eqref{eq:2.3} for $\phi$.  Formula \eqref{eq:9.88} and its
generalizations lead to the set of identities
\begin{equation}
  \Bigl\langle \frac{\delta S}{\delta\phi(x)}\ \phi(x_1)\cdots\phi(x_n)
    \Bigr\rangle = -i\sum_{i=1}^{n}\ \delta(x-x_i)\
    \langle \phi(x_1)\cdots\phi(x_n) \rangle .
  \label{eq:9.89}
\end{equation}
In this equation, the angle-brackets denote a time-ordered correlation
function in which derivatives on $\phi(x)$ are placed outside the
time-ordering symbol, as in Eq.~\eqref{eq:9.87}.  Relation \eqref{eq:9.89}
states that the classical Euler--Lagrange equations of the field $\phi$ are
obeyed for all Green's functions of $\phi$, up to contact terms arising from
the nontrivial commutation relations of field operators.  These quantum
equations of motion for Green's functions, including the proper contact terms,
are called \emph{Schwinger--Dyson equations}.

\subsection{Conservation Laws}

In classical field theory, Noether's theorem says that, to each symmetry of a
local Lagrangian, there corresponds a conserved current.  In Section~2.2 we
proved Noether's theorem by subjecting the Lagrangian to an infinitesimal
symmetry variation.  In the spirit of the above discussion of equations of
motion, we should find the quantum analogue of this theorem by subjecting the
functional integral to an infinitesimal change of variables along the symmetry
direction.

Again, it will be most instructive to begin with an example.  Let us consider
the theory of a free, complex-valued scalar field, with the Lagrangian
\begin{equation}
  \mathcal{L} = |\partial_\mu\phi|^2 - m^2|\phi|^2 .
  \label{eq:9.90}
\end{equation}
This Lagrangian is invariant under the transformation $\phi\to e^{i\alpha}\phi$.
The classical Noether current is
\begin{equation}
  j^\mu = i\bigl[ (\partial^\mu\phi^{*})\phi - \phi^{*}(\partial^\mu\phi) \bigr]
  = \frac{\partial\mathcal{L}}{\partial(\partial_\mu\phi)}\,\Delta\phi
    - \mathcal{J}^\mu ,
  \label{eq:9.96}
\end{equation}
which is conserved, $\partial_\mu j^\mu = 0$.  To find the quantum analogue,
we make the infinitesimal symmetry transformation
\begin{equation}
  \phi \to (1 + i\epsilon(x))\phi,\qquad
  \phi^{*} \to (1 - i\epsilon(x))\phi^{*},
\end{equation}
with $\epsilon(x)$ a function of spacetime.  Since the Lagrangian is invariant
for constant $\epsilon$, the variation of $\mathcal{L}$ is proportional to
$\partial_\mu\epsilon$.  Carrying through the steps leading up to
\eqref{eq:9.85}, we find the Schwinger--Dyson equation:
\begin{equation}
  \begin{aligned}
    \partial_\mu\ \bigl\langle j^\mu(x)\,\phi(x_1)\phi^{*}(x_2) \bigr\rangle
    &= -i\Bigl\langle \bigl( \Delta\phi(x_1)\,\delta(x-x_1) \bigr)
      \phi^{*}(x_2) + \phi(x_1)\bigl( \Delta\phi^{*}(x_2)\,
        \delta(x-x_2) \bigr) \Bigr\rangle .
  \end{aligned}
  \label{eq:9.98}
\end{equation}
A similar equation can be found for the correlator of $j^\mu$ with $n$ fields
$\phi(x)$.  These give the full set of Schwinger--Dyson equations associated
with the classical Noether theorem.

As an example of the use of this variational procedure to obtain the Noether
current, consider the symmetry of the Lagrangian with respect to spacetime
translations.  Under the transformation
\begin{equation}
  \phi_a \to \phi_a + a^\mu(x)\partial_\mu\phi_a ,
  \label{eq:9.99}
\end{equation}
the Lagrangian transforms as
\begin{equation}
  \mathcal{L} \to \mathcal{L} +
    \frac{\partial\mathcal{L}}{\partial(\partial_\nu\phi_a)}\,
      \partial_\mu a^\mu\,\partial_\nu\phi_a + a^\mu\partial_\mu\mathcal{L} .
\end{equation}
The variation of $\int d^4x\ \mathcal{L}$ with respect to $a^\mu$ then gives
rise to the conservation equation for the energy-momentum tensor
$\partial_\nu T^{\mu\nu} = 0$, with
\begin{equation}
  T^{\mu\nu} = \frac{\partial\mathcal{L}}{\partial(\partial_\mu\phi_a)}\,
    \partial^\nu\phi_a - g^{\mu\nu}\mathcal{L} ,
  \label{eq:9.100}
\end{equation}
in agreement with Eq.~\eqref{eq:2.17}.

The trick we have used in this section, that of considering a symmetry
transformation whose parameter is a function of spacetime, is reminiscent of a
technical feature of our earlier discussion introducing the Lagrangian of QED.
In Eq.~\eqref{eq:4.6}, we noted that the minimal coupling prescription for
coupling the photon to charged fields produces a Lagrangian invariant not only
under the global symmetry transformation with $\alpha$ constant, but also
under a transformation in which the symmetry parameter depends on $x$.  In
Chapter~15, we will draw these two ideas together in a general discussion of
field theories with local symmetries.

\subsection{The Ward-Takahashi Identity}

As a final application of the methods of this section, let us derive the
Schwinger--Dyson equations associated with the global symmetry of QED.
Consider making, in the QED functional integral, the change of variables
\begin{equation}
  \psi(x) \to (1 + ie\alpha(x))\psi(x),
  \label{eq:9.101}
\end{equation}
without the corresponding term in the transformation law for $A_\mu$ (which
would make the Lagrangian invariant under the transformation).  The QED
Lagrangian \eqref{eq:4.3} then transforms according to
\begin{equation}
  \mathcal{L} \to \mathcal{L} - j^\mu\partial_\mu\alpha ,
  \label{eq:9.102}
\end{equation}
where $j^\mu = \bar{\psi}\gamma^\mu\psi$ is the Dirac vector current.  The
transformation \eqref{eq:9.101} thus leads to the following identity for the
functional integral over two fermion fields:
\begin{equation}
  \begin{aligned}
    0 &= \int \mathcal{D}\bar{\psi}\,\mathcal{D}\psi\,\mathcal{D}A\
      e^{i\int d^4x\ \mathcal{L}}\ \Bigl[ -i\int d^4x\ \partial_\mu\alpha\,
      j^\mu(x)\ \psi(x_1)\bar{\psi}(x_2)\\
    &\quad + \bigl( ie\alpha(x_1)\psi(x_1) \bigr)\bar{\psi}(x_2)
      + \psi(x_1)\bigl( -ie\alpha(x_2)\bar{\psi}(x_2) \bigr) \Bigr] ,
  \end{aligned}
  \label{eq:9.103}
\end{equation}
with $j^\mu = \bar{\psi}\gamma^\mu\psi$.  As in our other examples, an analogous
equation holds for any number of fermion fields.

To understand the implications of this set of equations, consider first the
specific case \eqref{eq:9.103}.  Dividing this relation by $Z$, we find
\begin{equation}
  \begin{aligned}
    \langle 0|T\,\partial_\mu j^\mu(x)\,\psi(x_1)\bar{\psi}(x_2)|0\rangle
    &= -ie\,\delta(x-x_1)\ \langle 0|T\psi(x_1)\bar{\psi}(x_2)|0\rangle\\
    &\quad + ie\,\delta(x-x_2)\ \langle 0|T\psi(x_1)\bar{\psi}(x_2)|0\rangle .
  \end{aligned}
  \label{eq:9.104}
\end{equation}
To put this equation into a more familiar form, compute its Fourier transform
by integrating:
\begin{equation}
  \int d^4x\ e^{ik\cdot x}\ \int d^4x_1\,d^4x_2\ e^{-ip\cdot x_1}\
    e^{iq\cdot x_2}\ \cdots .
  \label{eq:9.105}
\end{equation}
Then the amplitudes in \eqref{eq:9.104} are converted to the amplitudes
$\mathcal{M}^\mu(k;p;q)$ and $\mathcal{M}_0(p;q)$ defined below \eqref{eq:7.67}
in our discussion of the Ward--Takahashi identity.  Indeed, \eqref{eq:9.104}
falls directly into the form
\begin{equation}
  -ik_\mu\mathcal{M}^\mu(k;p;q) = -ie\mathcal{M}_0(p;q-k)
    + ie\mathcal{M}_0(p+k;q) .
  \label{eq:9.106}
\end{equation}
This is exactly the Ward--Takahashi identity for two external fermions, which
we derived diagrammatically in Section~7.4.  It is not difficult to check that
the more general relations involving $n$ fermion fields lead to the general
Ward--Takahashi identity presented in \eqref{eq:7.68}.  Because of this
relation, the formula \eqref{eq:9.98} associated with the arbitrary symmetry
Eq.~(9.94) is usually also referred to as a Ward--Takahashi identity, the
one associated with the symmetry and its Noether current.

We have now arrived at a more general understanding of the terms on the
right-hand side of the Ward--Takahashi identity.  These are the contact terms
that we now expect to find when we convert classical equations of motion to
Schwinger--Dyson equations for quantum Green's functions.  The functional
integral formalism allows a simple and elegant derivation of these
quantum-mechanical terms.

\problems

\textbf{9.1} \emph{Scalar QED.}  This problem concerns the theory of a complex
scalar field $\phi$ interacting with the electromagnetic field $A^\mu$.  The
Lagrangian is
\begin{equation}
  \mathcal{L} = -\frac{1}{4}F^2_{\mu\nu} + (D_\mu\phi)^{*}(D^\mu\phi)
    - m^2\phi^{*}\phi ,
\end{equation}
where $D_\mu = \partial_\mu + ieA_\mu$ is the usual gauge-covariant
derivative.
\begin{enumerate}[label=(\alph*)]
  \item Use the functional method of Section~9.2 to show that the propagator
    of the complex scalar field is the same as that of a real field:
    \begin{equation}
      \phi(p)\phi^{*}(p') = \frac{i}{p^2-m^2+i\epsilon}\,
        (2\pi)^4\delta^{(4)}(p-p') .
    \end{equation}
    Also derive the Feynman rules for the interactions between photons and
    scalar particles; you should find
    \begin{equation}
      \text{(three-point vertex)} = 2ie(p_1+p_2)^\mu; \qquad
      \text{(four-point vertex)} = 2ie^2g^{\mu\nu} .
    \end{equation}
  \item Compute, to lowest order, the differential cross section for
    $e^{+}e^{-}\to\phi\phi^{*}$.  Ignore the electron mass (but not the scalar
    particle's mass), and average over the electron and positron
    polarizations.  Find the asymptotic angular dependence and total cross
    section.  Compare your results to the corresponding formulae for
    $e^{+}e^{-}\to\mu^{+}\mu^{-}$.
  \item Compute the contribution of the charged scalar to the photon vacuum
    polarization, using dimensional regularization.  Note that there are two
    diagrams.  To put the answer into the expected form,
    \begin{equation}
      \Pi^{\mu\nu}(q^2) = (g^{\mu\nu}q^2 - q^\mu q^\nu)\,\Pi(q^2) ,
    \end{equation}
    it is useful to add the two diagrams at the beginning, putting both terms
    over a common denominator before introducing a Feynman parameter.  Show
    that, for $-q^2 \gg m^2$, the charged boson contribution to $\Pi(q^2)$ is
    exactly $1/4$ that of a virtual electron-positron pair.
\end{enumerate}

\textbf{9.2} \emph{Quantum statistical mechanics.}
\begin{enumerate}[label=(\alph*)]
  \item Evaluate the quantum statistical partition function
    \begin{equation}
      Z = \mathrm{tr}\bigl[ e^{-\beta H} \bigr]
    \end{equation}
    (where $\beta = 1/k_BT$) using the strategy of Section~9.1 for evaluating
    the matrix elements of $e^{-iHt}$ in terms of functional integrals.  Show
    that one again finds a functional integral, over functions defined on a
    domain that is of length $\beta$ and periodically connected in the time
    direction.  Note that the Euclidean form of the Lagrangian appears in the
    weight.  Evaluate this integral for a simple harmonic oscillator,
    \begin{equation}
      L = \tfrac{1}{2}\dot{x}^2 + \tfrac{1}{2}\omega^2x^2 ,
    \end{equation}
    by introducing a Fourier decomposition of $x(t)$:
    \begin{equation}
      x(t) = \sum_n x_n\ \frac{1}{\sqrt{\beta}}\ e^{2\pi int/\beta} .
    \end{equation}
    The dependence of the result on $\beta$ is a bit subtle to obtain
    explicitly, since the measure for the integral over $x(t)$ depends on
    $\beta$ in any discretization.  However, the dependence on $\omega$ should
    be unambiguous.  Show that, up to a (possibly divergent and
    $\beta$-dependent) constant, the integral reproduces exactly the familiar
    expression for the quantum partition function of an oscillator.  [You may
    find the identity
    \begin{equation}
      \sinh z = z\cdot\prod_{n=1}^{\infty}\Bigl( 1 + \frac{z^2}{n^2\pi^2}
        \Bigr)
    \end{equation}
    useful.]
  \item Generalize this construction to field theory.  Show that the quantum
    statistical partition function for a free scalar field can be written in
    terms of a functional integral.  The value of this integral is given
    formally by
    \begin{equation}
      \bigl[ \det(-\partial^2+m^2) \bigr]^{-1/2} ,
    \end{equation}
    where the operator acts on functions on Euclidean space that are periodic
    in the time direction with periodicity $\beta$.  As before, the $\beta$
    dependence of this expression is difficult to compute directly.  However,
    the dependence on $m^2$ is unambiguous.  (More generally, one can usually
    evaluate the variation of a functional determinant with respect to any
    explicit parameter in the Lagrangian.)  Show that the determinant indeed
    reproduces the partition function for relativistic scalar particles.
  \item Now let $\psi(t)$, $\bar{\psi}(t)$ be two Grassmann-valued
    coordinates, and define a fermionic oscillator by writing the Lagrangian
    \begin{equation}
      L = \bar{\psi}i\partial_t\psi + \omega\bar{\psi}\psi .
    \end{equation}
    This Lagrangian corresponds to the Hamiltonian
    \begin{equation}
      H = \omega[\bar{\psi}, \psi] \qquad\text{with } \{\bar{\psi},\psi\}=1;
    \end{equation}
    that is, to a simple two-level system.  Evaluate the functional integral,
    assuming that the fermions obey antiperiodic boundary conditions:
    $\psi(t+\beta) = -\psi(t)$.  (Why is this reasonable?)  Show that the
    result reproduces the partition function of a quantum-mechanical two-level
    system, that is, of a quantum state with Fermi statistics.
  \item Define the partition function for the photon field as the
    gauge-invariant functional integral
    \begin{equation}
      Z = \int \mathcal{D}A\ \exp\Bigl\{ -\int d^4x_E\
        \tfrac{1}{4}(F_{\mu\nu})^2 \Bigr\}
    \end{equation}
    over vector fields $A^\mu$ that are periodic in the time direction with
    period $\beta$.  Apply the gauge-fixing procedure discussed in
    Section~9.4 (working, for example, in Feynman gauge).  Evaluate the
    functional determinants using the result of part (c) and show that the
    functional integral does give the correct quantum statistical result
    (including the correct counting of polarization states).
\end{enumerate}
